\documentclass{article}
\usepackage{amsmath}
\usepackage{amsthm}
\usepackage{fancyhdr}

\usepackage[backend=biber,sorting=none]{biblatex}
\addbibresource{mybib.bib}

\pagestyle{fancy}
\fancyhead[R]{ROUGH DRAFT}

\newtheorem{theorem}{Theorem}
\newtheorem{proposition}{Proposition}
\newtheorem{corollary}{Corollary}
\newtheorem{remark}{Remark}
\newtheorem{definition}{Definition}

\newcommand{\F}{\mathcal{F}}

\title{Maximal recoupling}
\author{Malcolm Sharpe}
\date{\today}

\begin{document}

\maketitle
% Apply fancy style to the first page
\thispagestyle{fancy}

\begin{abstract}
Given two random variables,
the maximal coupling algorithm samples from a new joint distribution that
maximizes agreement probability while preserving marginal distributions.
However, the structure of the original joint distribution is lost---in particular,
maximal subcouplings are typically destroyed.
We present a simple generalization,
the maximal recoupling algorithm,
which samples from a maximal coupling while preserving pre-existing maximal subcouplings.
We apply maximal recoupling to obtain a simultaneously maximal coupling of independent random variables.
\end{abstract}

\section*{Introduction}

\section*{Preliminaries}

Let $X$ and $Y$ be random variables taking values in a finite set $E$.

\begin{definition}
    A \textit{coupling} of $X$ and $Y$ is a joint distribution $\hat{P}$
    that preserves the marginal distributions.
    i.e. $\hat{P}_X = P_X$ and $\hat{P}_Y = P_Y$.
\end{definition}

\begin{definition}
    A coupling $\hat{P}$ is \textit{maximal} if $\hat{P}[X = Y] = 1 - D_{\text{TVD}}(P_X, P_Y)$.
\end{definition}

\subsection*{The maximal coupling algorithm}

Johnson \cite{johnson1998coupling} gave a rejection sampling algorithm for sampling from a maximal coupling.

\noindent\textbf{Given:} Sampling routines for \(P_X\) and \(P_Y\) and likelihood functions \(P_X\) and \(P_Y\).

\begin{enumerate}
    \item Generate \(x \sim P_X\) and \(u \sim U(0,1)\).
    \item If \(u < \frac{P_Y(x)}{P_X(x)}\), return \((x, x)\).
    \item Repeat:
    \begin{enumerate}
        \item Generate \(y \sim P_Y\) and \(u \sim U(0,1)\).
        \item If \(u > \frac{P_X(y)}{P_Y(y)}\), return \((x, y)\).
    \end{enumerate}
\end{enumerate}

\section*{The maximal recoupling algorithm}

\subsection*{Pseudocode}

\noindent\textbf{Given:} A joint distribution \(P_{XY}\) providing sampling routines for \(P_X\) and \(P_{Y|X=x}\) and likelihood functions \(P_X\) and \(P_Y\).

\begin{enumerate}
    \item Generate \(x \sim P_X\) and \(u \sim U(0,1)\).
    \item If \(u < \frac{P_Y(x)}{P_X(x)}\), return \((x, x)\).
    \item Set \(a \gets x\).
    \item Repeat:
    \begin{enumerate}
        \item Generate \(y \sim P_{Y|X=a}\) and \(u \sim U(0,1)\).
        \item If \(u > \frac{P_X(y)}{P_Y(y)}\), return \((x, y)\).
        \item Set \(a \gets y\).
    \end{enumerate}
\end{enumerate}

\begin{remark}
    When $X$ and $Y$ are independent, this algorithm is equivalent to the rejection sampling
    algorithm for maximal coupling.
\end{remark}

\subsection*{Analysis}

Let \(\hat{P}_{XY}\) be the joint distribution arising from the maximal recoupling algorithm.
The first few properties also hold for the maximal coupling algorithm.

\begin{proposition}
    \(\hat{P}_{XY}\) is a coupling. i.e. $\hat{P}_X = P_X$ and $\hat{P}_Y = P_Y$.
\end{proposition}

\begin{proposition}
    \(\hat{P}_{XY}\) is maximal. i.e. For all \(e \in E\),
    \[
        \hat{P}_{XY}(e,e) = \min \{ P_X(e), P_Y(e) \}.
    \]
\end{proposition}

\begin{proposition}
    Assume a series of draws from $\hat{P}_{Y|X}$ obeys,
    for all $x \in E$,
    \[
        E[\text{number of draws from \(\hat{P}_{Y|X=x}\)}] \le P_X(x).
    \]
    Then, for all $x \in E$,
    \[
        E[\text{number of draws from \(P_{Y|X=x}\)}] \le P_X(x).
    \]
\end{proposition}

\begin{corollary}
    During one draw from $\hat{P}_{XY}$,
    \[
        E[\text{number of draws from \(P_{Y|X}\)}] \le 1.
    \]
\end{corollary}

The next property is what distinguishes maximal recoupling from maximal coupling.

\begin{proposition}
    Let $E'$ be a finite set and $f : E \to E'$. Then
    \[
    D_{\text{TVD}}(P_{f(X)f(Y)}, \hat{P}_{f(X)f(Y)})
    \le 2 \left(\hat{P}[f(X) = f(Y)] - P[f(X) = f(Y)]\right).
    \]
\end{proposition}

\begin{corollary}
    Let $E'$ be a finite set and $f : E \to E'$ such that
    $P_{f(X)f(Y)}$ is a maximal coupling of $f(X)$ and $f(Y)$.
    Then $\hat{P}_{f(X)f(Y)} = P_{f(X)f(Y)}$.
\end{corollary}

We can express $\hat{P}$ in closed form.

\begin{proposition}
    For all $x, y \in E$:
    \begin{itemize}
        \item If $x = y$, then $\hat{P}_{XY}(x,y) = \min \{ P_X(x), P_Y(y) \}$.
        \item If $x \neq y$:
        \begin{itemize}
            \item If $P_X(x) \leq P_Y(x)$ or $P_Y(y) \leq P_X(y)$, then $\hat{P}_{XY}(x,y) = 0$.
            \item Otherwise,
\begin{align*}
&\hat{P}_{XY}(x,y) = \\
&\quad\max \{ 0, P_X(x) - P_Y(x) \} \\
&\quad\max \{ 0, P_Y(y) - P_X(y) \} \\
&\quad\left( P_X(x) P_Y(y) \right)^{-1/2} \\
&\quad\left( (I - A)^+ \right)_{xy}
\end{align*}
                where, for all $i,j \in E$,
\[
A_{ij} =
\begin{cases}
    \frac{P_{XY}(i,j)}{\sqrt{P_{\text{max}}(i) P_{\text{max}}(j)}} & \text{if } P_{\text{max}}(i) P_{\text{max}}(j) > 0, \\
    0 & \text{otherwise},
\end{cases}
\]
                where, for all $i \in E$,
\[
P_{\text{max}}(i) = \max \{ P_X(i), P_Y(i) \}.
\]
        \end{itemize}
    \end{itemize}
\end{proposition}

\section*{Simultaneously maximal coupling}

For $n \ge 1$, let \( \mathcal{F} \) be a family of subsets of \([n]\).
In this section, \( X = (X_1, \dots, X_n) \), \( Y = (Y_1, \dots, Y_n) \).

\begin{definition}
$\F$ is simultaneously maximally coupled if, for all $S \in \F$,
$X_S$ and $Y_S$ are maximally coupled. i.e.
\[
    P[X_S = Y_S] = 1 - D_{\text{TVD}}(P_{X_S}, P_{Y_S}).
\]
\end{definition}

This first theorem is well-known in the Markov chain literature.

\begin{theorem}
The following are equivalent:
\begin{itemize}
    \item For every choice of marginal distributions for $X$ and $Y$,
    $\F$ can be simultaneously maximally coupled.
    \item For every $S,T \in \F$, either $S \subseteq T$ or $T \subseteq S$.
\end{itemize}
\end{theorem}

The independent case is novel.

\begin{theorem}
\raggedright
The following are equivalent:
\begin{itemize}
    \item For every choice of marginal distributions for $X$ and $Y$
    such that $X_1, \dots, X_n$ are independent
    and $Y_1, \dots, Y_n$ are independent,
    $\F$ can be simultaneously maximally coupled.
    \item $\F$ is laminar. i.e. For every $S,T \in \F$, either $S \subseteq T$, $T \subseteq S$, or $S \cap T = \emptyset$.
\end{itemize}
\end{theorem}

These two theorems can be generalized by Bayesian network structure.

\begin{theorem}
\raggedright
Let $G$ be a directed graph with vertex set $[n]$. The following are equivalent:
\begin{itemize}
    \item For every choice of marginal distributions for $X$ and $Y$
    such that $X_1, \dots, X_n$ and
    $Y_1, \dots, Y_n$ are both Bayesian networks with respect to $G$,
    $\F$ can be simultaneously maximally coupled.
    \item $\F$ is laminar and, for every disjoint $S,T \in F$, $S$ and $T$ are d-separated by $\emptyset$.
\end{itemize}
\end{theorem}

For computational purposes, this corollary is more convenient.

\begin{corollary}
Let $G$ be a directed graph with vertex set $[n]$.
Assume $\F$ is laminar and ancestral.
i.e. Every $S \in \F$ contains all ancestors of each of its elements.

Then, for every choice of marginal distributions for $X$ and $Y$
such that $X_1, \dots, X_n$ and
$Y_1, \dots, Y_n$ are both Bayesian networks with respect to $G$,
$\F$ can be simultaneously maximally coupled.
\end{corollary}

\section*{Related Work}

\textbf{Maximal coupling.}
Maximal coupling is generally credited to Doeblin \cite{lindvall1991w}.
The rejection sampling algorithm for maximal coupling was introduced by Johnson \cite{johnson1998coupling}.
The algorithm was used implicitly for speculative decoding of auto-regressive large language models \cite{leviathan2023fast},
which inspired our work;
the connection was later made explicit by Sun et al. \cite{sun2024spectr}.

\textbf{Simultaneously maximal coupling.}
Sun et al. \cite{sun2024blockverificationacceleratesspeculative_v2}
thoroughly investigated simultaneously maximal coupling for speculative decoding
of auto-regres\-sive large language models,
under the name Greedy Block Verification.
The key ideas can be previously found in the Markov chain literature.

\textbf{Coupling with more than two random variables.}
Another interesting direction for generalizing maximal coupling
is with three or more random variables \cite{angel2021pairwiseoptimalcouplingmultiple}.

\printbibliography

\end{document}
